% Transcription — fonds Alexandre Grothendieck, shelfmark 71, batch 1 (pages 1-20).
% Established from the facsimile archives/batches/71/batch-01.pdf (a mirror of
% https://grothendieck.umontpellier.fr/71.pdf), per the skill
% transcribe-grothendieck. The batch file carries no cover sheet: PDF page k is
% archivists' page k; the pencilled numbers bottom-left (« 1 » ... « 20 »)
% agree.
%
% Pass: Opus 5.5 (claude-opus-5-5), 2026-09-24 — first pass, unchecked against the pages by a human.
%
% Eleven of the twenty pages carry a \page marker. Absent: pages 1, 4, 15, 17
% and 19, blank cream sheets (separators, bearing only the pencilled number);
% pages 3 and 14, typed USTL « avis de soutenance » dated Montpellier, 19 avril
% 1977 (theses in geology and chemistry), with nothing in his hand; page 8, a
% printed Crédit Lyonnais leaflet dated « Janvier 1978 », upside down, with
% nothing in his hand. Those two dated versos are consistent with the
% archivists' « [à partir de 1976] » and say nothing more.
%
% Whose hand. Pages 2, 5-7, 9-13, 16 and 18 are his working notes, transcribed
% in full. Page 20 is not: it is page « I » of un mémoire d'une autre main,
% signé et daté du 15 juillet 1976 (blue ink, a rounded hand), whose last
% leaves XIV-XIX are pages 41-47 (batch 3). By the user's decision of
% 2026-09-24, and following folder 100's rule for other people's text
% (transcripts/100/batch-02.fr.tex), the mémoire is summarised in one French
% \note per page, not recomposed; the author is not named. The marks on it
% are given as \marginal, each tied to the words they answer: those in black
% ink, in the hand of his working leaves (« or. ? », « et non vide », « de
% R^2 », the margin dots, « faux »), and the pencil marks of an unidentified
% hand (« plan », the dashed box round « Remarque 1 »).
%
% What the batch is. Three strands of the Jordan-curve question, none
% paginated by him, the order being the archivists'.
%   2, 13  Theta-graphs: three arcs Γ'_1, Γ'_2, Γ'_3 joining two points
%          ∂ = {x, y} in a plane E, the three curves Γ_i = Γ'_j ∪ Γ'_k and
%          the regions V_1, V_2, V_3; page 13 states the equivalence
%          a) the three regions near x are distinct / a') E - Γ has three
%          regions / b) Γ_1, Γ_2, Γ_3 are Jordan curves / c) Γ_1, Γ_2 are,
%          and proves c) => a).
%   9, 11, 7, 5  A numbered sequence 1)-6) (his numbering, pages out of
%          order): pairs of open sets each the complement of the other's
%          closure (1); a closed K with two relatively open U_K, U'_K sharing
%          their frontier C (3); pseudo-partitions U, U', W, B, B', C of X
%          (4); gluing X_i = X - U_i back along the B_i (5); the space 𝒳
%          obtained by gluing, étale over X, and an « immersion » lemma (6).
%   6, 12  A cancelled computation: a poset E, the realisations |E_e|, a
%          map φ : M(E) -> I to an ordered set, the filtration X_i and the
%          complexes Σ_d(𝒳) indexed by flags d ∈ Drap(I), with
%          π_0 Σ_d = D(d) ≅ E(d); on page 12 also the pair (K, L) of
%          simplicial complexes and Σ(|K|, |L|) as a colimit of simplices.
%   10     Unrelated algebra in his hand: an automorphism t of a group on
%          λ_0, λ_1, λ_2 with t^3 = int(λ_0) (upside down on the sheet), and a
%          pencil sketch with G ×_H E and B_H -> B_G. Transcribed, condensed.
%   16     Theta-graphs again with an index set I (B ≅ I × {0,1},
%          U_i, V_i), then three equivalent conditions for a simple closed
%          curve C in a connected surface to separate it.
%   18     Proposition: two J-polygons Γ, Γ' meeting in a polygonal arc I
%          are either exterior to each other or nested; the J-polygon
%          Γ'' = (Γ ∪ Γ') - (I - ∂I) and its interior and exterior in each case.
%   20     Page I of the mémoire (another hand), summarised: Théorèmes 1
%          and 2, the definitions of § 1, Remarque 1; his marks as marginals.
%
% Notation, fixed once. His dotted letter (U̇, K̇, Ẇ) is the frontier and is
% set \dot{}; his ∁ (a tall C) is \complement; « Int », « Ext » are
% \operatorname. His script X (𝒳) is \mathcal{X}. Γ', Γ'' and Γ° are as drawn.
% His own interlinear insertions on his working leaves are integrated
% silently into the line; his cancellations are \struck{}.
%
% The dating is the folder's own; the group « Géométrie et topologie
% combinatoire » (69-89) holds it.
\documentclass[11pt,a4paper]{article}
% Le préambule commun à toutes les transcriptions du fonds Grothendieck.
%
% Il tient en une idée : **l'incertitude se compose, elle ne se résout pas.**
% Une transcription qui lisse un mot douteux a détruit la seule chose pour
% laquelle on la faisait — un lecteur doit pouvoir savoir, à chaque endroit,
% ce qui était sur le feuillet et ce qui a été ajouté. Les six macros qui
% suivent sont donc l'appareil critique, et non de la décoration.
%
% Les mêmes commandes sont reconnues par `scripts/render.mjs`, qui produit la
% vue de lecture du site. Une macro ajoutée ici doit l'être aussi là-bas,
% faute de quoi le rendu échouera bruyamment — ce qui est voulu : un rendu
% silencieusement faux est pire qu'un rendu absent.

\ProvidesPackage{grothendieck}[2026/08/08 Transcriptions du fonds Grothendieck]

\usepackage{amsmath,amssymb,amsthm}
\usepackage{mathrsfs}
\usepackage{stmaryrd}
% Les diagrammes commutatifs sont partout dans ce fonds : c'est la langue
% même de la géométrie algébrique de Grothendieck, et une transcription qui
% les rendrait en prose ne transcrirait rien.
\usepackage{tikz-cd}
% Français par défaut — c'est la langue des feuillets — mais l'anglais est
% chargé aussi : la traduction et le résumé sont des documents anglais, et la
% typographie française (espaces avant les deux-points) y serait une faute.
% Ces documents basculent par \selectlanguage{english} en tête de corps.
\usepackage[english,french]{babel}
\usepackage{fontspec}
\usepackage{geometry}
\geometry{a4paper,margin=25mm,marginparwidth=28mm}
\usepackage{marginnote}
\usepackage{xcolor}
% ulem, not soul. `soul`'s \st and \ul are built on a character-by-character
% scan that XeLaTeX's font machinery breaks: the document fails with an
% undefined control sequence rather than degrading. ulem does the same two jobs
% and is Unicode-engine safe. `normalem` stops it hijacking \emph, which the
% transcriptions use for Grothendieck's own emphasis.
\usepackage[normalem]{ulem}
\usepackage{hyperref}
\hypersetup{colorlinks=true,linkcolor=black,urlcolor=black,pdfborder={0 0 0}}

% --- Métadonnées du lot -----------------------------------------------------
% Reprises telles quelles par le script de rendu, qui les affiche en tête de
% la vue de lecture. Elles fixent ce que le fichier prétend couvrir : sans
% elles, une transcription flotte sans point d'attache dans un fonds de seize
% mille feuillets.

\newcommand{\folder}[1]{\gdef\grothFolder{#1}}
\newcommand{\batch}[1]{\gdef\grothBatch{#1}}
\newcommand{\leaves}[2]{\gdef\grothFirst{#1}\gdef\grothLast{#2}}
\newcommand{\foldertitle}[1]{\gdef\grothFoldertitle{#1}}
\gdef\grothFolder{?}\gdef\grothBatch{?}\gdef\grothFirst{?}\gdef\grothLast{?}\gdef\grothFoldertitle{}

% --- Le résumé --------------------------------------------------------------
%
% Il ouvre la lecture modernisée, sous le titre « Résumé », et il est composé
% en petit. Ce n'est pas un ornement : c'est le seul passage du document qui
% ne soit pas des mathématiques, et un lecteur doit pouvoir le reconnaître
% avant de l'avoir lu — puis le sauter s'il connaît déjà le sujet, ou s'y
% arrêter s'il ne le connaît pas. Le filet qui le suit marque la reprise du
% corps.

\newenvironment{resume}
  {\small\setlength{\parskip}{0.45em}}
  {\par\medskip\hrule\medskip}

% --- Mots-clés ---------------------------------------------------------------
%
% En anglais, en clôture du résumé de la lecture modernisée : le vocabulaire
% moderne sous lequel chercher ce que les feuillets construisent. Le manifeste
% du site (`npm run manifest`) les extrait pour étiqueter la cote entière —
% cette ligne est la source unique des tags, il n'y a pas de fichier à côté.
\newcommand{\keywords}[1]{\par\smallskip\noindent
  {\footnotesize\textsc{Keywords} --- \textit{#1}}\par}

% --- L'appareil critique ----------------------------------------------------

% Le numéro de feuillet, en marge. C'est l'ancre qui permet de régler un
% désaccord sur une formule en regardant *un* feuillet plutôt qu'une chemise
% de six cents. Le site s'en sert aussi pour tourner les pages du fac-similé
% quand on fait défiler la transcription.
\newcommand{\leaf}[1]{%
  \marginnote{\footnotesize\color{gray!70}\textsf{#1}}%
  \label{leaf:#1}\ignorespaces}

% Illisible : on ne devine pas. Un mot inventé qui se lit comme les autres est
% le pire résultat possible.
\newcommand{\ill}{\textcolor{red!60!black}{[\,\ldots\,]}}

% Lecture douteuse : le mot est donné, et signalé comme incertain.
\newcommand{\uncertain}[1]{\uline{#1}}

% Ajout de l'éditeur — une lettre manquante, un mot sous-entendu.
\newcommand{\add}[1]{\textcolor{blue!50!black}{[#1]}}

% Biffé par Grothendieck lui-même. Ce qu'il a rayé fait partie du document :
% une rature dit souvent où la pensée a changé de direction.
\newcommand{\struck}[1]{\sout{#1}}

% Note du transcripteur, sur l'état matériel du feuillet.
\newcommand{\note}[1]{\par\smallskip{\small\color{gray!60!black}\textit{[#1]}}\par\smallskip}

% Note marginale de Grothendieck, distincte du corps du texte.
\newcommand{\marginal}[1]{\par\smallskip\noindent\hspace{1em}%
  {\small\color{gray!60!black}#1}\par\smallskip}

% --- Titre ------------------------------------------------------------------

\AtBeginDocument{%
  \begin{center}
    {\large\bfseries Fonds Alexandre Grothendieck --- cote n\textsuperscript{o}~\grothFolder}\\[2pt]
    {\itshape\grothFoldertitle}\\[4pt]
    {\small feuillets \grothFirst--\grothLast\ (lot \grothBatch)}\\[2pt]
    {\footnotesize\color{gray!60!black}Transcription établie d'après le fac-similé numérisé
     par l'Université de Montpellier.\\
     Les lectures incertaines sont soulignées, les illisibles notées [\,\ldots\,],
     les ajouts entre crochets.}
  \end{center}
  \vspace{1em}}

\endinput


\folder{71}
\batch{1}
\pages{1}{20}
\dating{[à partir de 1976]}
\watermark{Édition de démonstration}
\foldertitle{Théorème de Jordan : notes manuscrites (1976, s.d.).}

\begin{document}

\page{2}
\note{en haut à gauche, trois colonnes ; à droite d'un trait vertical, les
bords}

\[
\begin{array}{ccc}
V_1 & V_2 & V_3 \\
\Gamma''_1 & \Gamma''_2 & \Gamma''_3 \\
\Vert & \Vert & \Vert \\
\Gamma'_1 - \partial & \Gamma'_2 - \partial & \Gamma'_3 - \partial
\end{array}
\qquad \partial = \{x, y\}
\]

$V_1$ \uncertain{bords} $\Gamma'_2, \Gamma'_3$ ; $\Gamma'_1$ bord
\uncertain{pour} $V_2, V_3$

$V_2$ \uncertain{bords} $\Gamma'_3, \Gamma'_1$ ; $\Gamma'_2$ — $V_3, V_1$

$V_3$ \uncertain{bords} $\Gamma'_1, \Gamma'_2$ ; $\Gamma'_3$ — $V_1, V_2$

\[
\begin{aligned}
\Gamma_1 &= \Gamma_2 \cup \Gamma_3 = \Gamma'_2 \cup \Gamma'_3 = \Gamma''_1 \cup \Gamma''_2 \cup \partial \\
\Gamma_2 &= \Gamma_3 \cup \Gamma_1 = \Gamma'_3 \cup \Gamma'_1 = \Gamma''_3 \cup \Gamma''_1 \cup \partial \\
\Gamma_3 &= \Gamma_1 \cup \Gamma_2 = \Gamma'_1 \cup \Gamma'_2 = \Gamma''_1 \cup \Gamma''_2 \cup \partial
\end{aligned}
\]
\note{ainsi sur la page, où les deux premiers membres portent des $\Gamma$
sans accent ; page 13 écrit $\Gamma_1 = \Gamma'_2 \cup \Gamma'_3$. Les
indices du dernier membre de la première ligne sont d'une lecture incertaine}

\[
\Gamma^{*}_3 = \Gamma''_1 \cup \Gamma''_2 \cup \partial
\quad\Big|\quad
U_3 = V_1 \cup V_2 \cup \Gamma''_3
\quad\Big|\quad
V_3
\]

$\dot{U}_3 = \Gamma_3$ ; $U_3 \cup$

$U, C, V$

\note{à gauche, un croquis : deux cercles sécants, hachurés obliquement, un
petit cercle autour d'un des points d'intersection ; non redessiné}

\page{5}
$\mathcal{X}$ \uncertain{peut aussi s'identifier à} la
\uncertain{somme amalgamée} des $\overline{W}$ et des
$\coprod_{i \in I} \overline{U_i}$ (via
$\coprod_i B_i \hookrightarrow \overline{W}$,
$\coprod_i B_i \to \coprod_i \overline{U_i}$) :

\struck{L'application naturelle}

\struck{$\mathcal{X} \to X$}

\[
\mathcal{X} \simeq \overline{W} \amalg_{(\coprod_i B_i)} \Bigl(\coprod_i \overline{U_i}\Bigr)
\simeq \coprod_{i\,(W)} X_i \qquad X_i \simeq \mathcal{X}_i
\]
\note{sous $\mathcal{X}_i$ : « $\Vert$ $\operatorname{Im} X_i$ », relié par
une flèche à « \uncertain{ouvert} de $\mathcal{X}$ »}

$\mathcal{X} \to X$ est étale en les pts de
$\mathcal{X} - \coprod_i C_i$, i.e.\ une \ill{} \uncertain{fermée} en les
pts de $\coprod C_i$.

\note{à gauche, une figure : trois régions arrondies accolées autour d'un
nœud à trois branches épaissies, marquées $B_1, B_2, B_3$ aux extrémités et
$C_1, C_2, C_3$ près du centre ; sous la figure, « $K$ » ; non redessinée}

6) \uncertain{Lemme} d'immersion de $\mathcal{X}$ \ill{}

\uncertain{remplaçant} au voisinage de $K$ par $K^{*} \subset K$
($U_i, U'_i$ par $U^{*}_i = U_i \cap K^{*}$, $U'^{*}_i = U'_i \cap K^{*}$),
\uncertain{on suppose} que $\dot{K}^{*} \subset U_i$, et
\struck{$\dot{K} \subset \bigcup U_i$, et} \struck{$(K \setminus K')$}
$= V_i$,
$C_i \subset \overline{U^{*}_i} \cap \overline{U'^{*}_i}$ (il suffit que
\struck{\ill{}} $\operatorname{int} K' \supset \bigcup C_i$.

— vérifier si dans les développements précédents, il n'y avait pas lieu
d'exiger $\bigcup C_i \subset \operatorname{int} K$ i.e.\
$C_i \cap \dot{K} = \emptyset$ $\forall i \in I$…

\marginal{vérifier \ill{} \ill{}…}

\page{6}
\note{page barrée d'un long trait oblique d'un bout à l'autre, et d'une
croix dans le quart inférieur droit}

\[
\begin{array}{c}
e \in E \\
\cup \\
M(e) = M(E_e) \subset E_0 = M(E) \\
\big\downarrow{\scriptstyle \varphi} \\
I \ \text{ordonné}
\end{array}
\qquad
\begin{cases}
|E| = X \supset |E_e| \\
X_i = \bigcup_{\varphi(M(e)) \leq i} |E_e|
\end{cases}
\qquad
E_e = \{x \in E \mid x \leq e\}
\]

$e_0 \in E_0 \mapsto X_{e_0} = \bigcup_{M(e) \leq e_0} |E_e|$

$(i_0 > \cdots > i_k) = d \in \operatorname{Drap}(I)$, $d \subset I$

$\Sigma_d(\mathcal{X})$ ;
$\Sigma_{\{i\}}(X) = \Sigma(X_i, \partial X_i)$

\[
\begin{array}{ccccccc}
\Sigma_{\{i_0, \dots, i_{k-1}\}}(\mathcal{X}) & \supset & \partial\Sigma_{\{i_0, \dots, i_{k-1}\}} & \supset & \Delta_{i_0 \dots i_k}(\mathcal{X}) & \supset & \partial\Delta_{i_0 \dots i_k} \\
\big\downarrow & & \big\downarrow & & \big\downarrow & & \big\downarrow \\
\Sigma_{i_{k-1}}(X)_{\ill{}} & & & & & & \\
\big\downarrow & & & & & & \\
X_{i_{k-1}} & \supset & \partial X_{i_{k-1}} & \longleftarrow & X_{i_k} & \supset & \partial X_{i_k}
\end{array}
\]
\note{les trois flèches verticales de droite descendent de la première
ligne à la dernière ; le sens de la flèche entre $\partial X_{i_{k-1}}$ et
$X_{i_k}$ est incertain}

\[
\Sigma_{\{i_0, \dots, i_k\}}(\mathcal{X}) = \Sigma(\Delta_{i_0 \dots i_k}, \partial\Delta_{i_0 \dots i_k})
\]

$\Sigma_{\{i_0, \dots, i_k\}} \to \Sigma_{d'}(\mathcal{X})$ pour
$d' \subsetneq d$ ; en colonne,
$\partial X_{i_k} \subset X_{i_k} \to X_{i_{k-1}}$.

$\Sigma_d \to \Sigma_{\tilde{d}'} \to \Sigma_{d'}$

\begin{tikzcd}
\Sigma_{d \setminus \{i_k\}} \arrow[r] \arrow[d] & X_{i_{k-1}} \arrow[d, hook] \\
\Sigma_{d' \setminus \{i_k\}} \arrow[r] & X_{i_\alpha}
\end{tikzcd}

\note{sous $\Sigma_{d' \setminus \{i_k\}}$, une flèche surchargée et biffée
vers $\Sigma_{d'}$ ; après $\Sigma_{d' \setminus \{i_k\}}$, un indice biffé}

$\pi_0(\Sigma_d) = D(d)$

\page{7}
5) Soit $K$ un fermé de l'espace $X$, $(U_i)_{i \in I}$ une famille
d'ouverts relatifs de $K$, mutuellement disjoints, tels que
$\dot{K} \subset \bigcup_i U_i$. Soit $B_i = \dot{K} \cap U_i$, de
sorte que $\dot{K} = \coprod B_i$. (\struck{Supposons} Pour tt $i$, soit
$C_i = \overline{U_i} \setminus U_i$ la frontière relative de $U_i$ dans
$K$, supposons que $C_i \subset (K - \overline{U_i})^{-}$ cf.\ …). Posons
\struck{\ill{}} $X_i = X \setminus U_i$. Alors
\note{l'indice de $U_i$ est surchargé à chaque occurrence de la page ; un
accent (\,$U'_i$\,) est lisible par endroits et n'est pas retenu}

\[
\begin{cases}
X_i \text{ est loc.\ fermé} \\
\operatorname{Int} X_i = X - \overline{U_i} = X_i \setminus C_i \\
\overline{X_i} = X_i \cup B'_i \qquad (B'_i = \dot{K} \setminus B_i)
\end{cases}
\]
\marginal{\struck{i.e.\ $C_i$}}

$\operatorname{Int} X_i$ et $U_i \setminus C_i$ sont \struck{deux} ouverts
de $X$ dont chacun est le complémentaire de l'adhérence de l'autre.

Soit $W = X \setminus K$, \struck{\ill{}} donc
$\dot{W} = \overline{W} \cap K = \dot{K} = \coprod B_i$, l'inclusion
$W \subset X_i$ est une immersion ouverte et
$\dot{W} \cap X_i = \dot{K} \cap X_i = B_i$, \struck{\ill{}} i.e.\
\struck{\ill{}} l'adh.\ de $W$ dans $X_i$ est $W \cup B_i$, dans la
\struck{bord} frontière de $W$ dans $X_i$ est $B_i$. On a
\[
X_i = \bigl(\overline{W} \cap X_i = W \cup B_i\bigr) \amalg_{B_i} \overline{U_i}
\]
\struck{\ill{}}

Alors la somme amalgamée des \ill{} \ill{} \ill{} $W$ est \ill{}
\note{la dernière ligne, au bord inférieur du feuillet, n'est lue qu'en
partie}

\page{9}
1) $X$ espace topologique, $U, V$ ouverts, conditions équivalentes
\begin{enumerate}
\item[a)] $U = \complement\overline{V}$, $V = \complement\overline{U}$
\item[b)] $U \cap V = \emptyset$,
$\complement(U \cup V) \subset \overline{U} \cap \overline{V}$ (d'où
égalité)
\item[c)] $U \cap V = \emptyset$, $X = \overline{U \cup V}$,
$\dot{U} = \dot{V}$
\item[d)] $\overline{U} \cap \overline{V} = \dot{U} = \dot{V}$.
\end{enumerate}

$U$ et $V$ se déterminent mutuellement. Pour $U$ fixé, il existe $V$ ayant
les propriétés précédentes ssi $V = \complement\overline{U}$ est dense dans
$\complement U$.

3) $X$ espace topologique, $K \subset X$ fermé, $U_K$ et $U'_K$ deux ouverts
relatifs de $K$ tels que
$\overline{U_K} \setminus U_K = \overline{U'_K} \setminus U'_K$,
$\overline{U_K} = K \setminus U'_K$, $\overline{U'_K} = K \setminus U_K$,
$C = \overline{U_K} \cap \overline{U'_K}$ $=$ frontière commune (fermé dans
$X$) de $U_K$ et $U'_K$ dans $K$ \struck{\ill{}}, soit
$B = U_K \cap \dot{K}$, \struck{donc} donc $B$ ouvert dans $\dot{K}$ et
fermé dans $U_K$ donc \uncertain{relativement} fermé dans $X$. Alors on a,
\uncertain{puisque} $B \cap C = \emptyset$.

\[
\begin{cases}
U = U_K \setminus B = \operatorname{Int}(U_K) \\
V = X \setminus \overline{U_K} = (X \setminus U_K) \setminus C = \operatorname{Int}(X_1),
\qquad X_1 = X \setminus U_K \\
U = \complement\overline{V}, \quad V = \complement\overline{U}, \quad
\dot{U} = \dot{V} = \overline{U} \cap \overline{V} = \complement(U \cup V) = B \cup C \\
\struck{U \cap K = U \subset K,\ V \cap K = V_K}\quad V \subset X_1 \subset \overline{V}
\end{cases}
\]

\emph{Cor.} \struck{Si $B \cap C = \emptyset$ i.e.\ \ill{}
$U_K \setminus U$. Supposons, Pour $B$} Pour que $X_1 = \complement U_K$
soit loc.\ fermé, il f.\ et s.\ que $B = U_K \cap \dot{K}$ soit
\struck{fermé dans $X$, i.e.\ $B$ soit ouvert et fermé} dans
$B \cup C$ \ill{} dans $X$, \ill{} on suppose que \ill{} soit une somme
\ill{} topologique \struck{de $\dot{K}$, et posons $\dot{K} = B' \amalg B$}
\struck{\ill{} $\dot{K}$. Posant alors $\dot{K} = B' \amalg B$} qu'il soit
fermé dans $B \cup C$, on suppose que $\overline{B} \cap C = \emptyset$.
[Il suffit pour ceci que $\dot{K} \cap C = \emptyset$ i.e.\ que
$U_K \cup V_K \supset \dot{K}$] \struck{i.e.\ $U_K$ et $V_K$ \ill{}}
\note{fin de page très raturée ; seules les formules et les ratures
lisibles sont données}

\page{10}
\note{feuillet plié, sans rapport visible avec le reste du dossier ; la
moitié gauche est écrite tête-bêche, la moitié droite porte un croquis au
crayon. Transcription condensée}

$t^2(\lambda_0) = \lambda_0$, ;
$t^2(\lambda_1) = t(\uncertain{q_\nu})\, t(\lambda_\nu)\, t(\uncertain{q_\nu})^{-1}$
\note{sous $t(\lambda_\nu)$, une accolade « $\lambda_1$ »}

$\lambda_0 \ \lambda_1 \ \lambda_2$ ;
$t^3 = \operatorname{int}(\lambda_0)$

\[
\begin{aligned}
t(\lambda_0) &= \lambda_0 \\
t(\lambda_1) &= \lambda_1 \lambda_2 \lambda_1^{-1} = \lambda_0^{-1} \lambda_2 \lambda_0 \\
t(\lambda_2) &= \lambda_1
\end{aligned}
\]

\[
\begin{aligned}
t^3(\lambda_0) &= \lambda_0 \\
t^3(\lambda_1) &= (\lambda_1 \lambda_2 \lambda_1^{-1})\, \lambda_1\, (\lambda_1 \lambda_2^{-1} \lambda_1^{-1})
= \lambda_1 \lambda_2 \lambda_1 \lambda_2^{-1} \lambda_1^{-1} \\
&= \operatorname{int}(\lambda_1 \lambda_2)\, \lambda_1 = \operatorname{int}(\lambda_0^{-1})\, \lambda_1 \\
t^3(\lambda_2) &= t(\lambda_1) = \lambda_1 \lambda_2 \lambda_1^{-1} = \struck{\operatorname{int}(\lambda_1)} \\
&= (\lambda_1 \lambda_2)\, \lambda_2\, (\lambda_2^{-1} \lambda_1^{-1})
= \operatorname{int}(\lambda_1 \lambda_2)\, \lambda_2\ \operatorname{int}\ill{}
\end{aligned}
\]
\note{dans la seconde ligne, deux facteurs $\lambda_1^{-1}$ et $\lambda_1$
sont barrés d'un trait oblique, simplifiés ; le texte porte
« $\operatorname{int}$ » en abrégé}

\note{à droite, au crayon, pâle : un diagramme
$G/H \leftarrow E = G \times_H E \leftarrow G \times_H E$, avec des flèches
verticales vers $E$ et $E'$, puis $B_H \to B_G$ (flèche $f$),
$f^{*}(E) \to E$, $f^{*} f_{!}(\ill{}) \to \ill{}$ et « $\tilde{E} \simeq$ » ;
lecture incertaine dans le détail}

\page{11}
3) Supposons $\dot{K} \cap C = \emptyset$, et posons \struck{\ill{}}
\begin{align*}
B' &= U'_K \cap \dot{K} \\
U' &= U'_K \setminus B' = \operatorname{Int}(U'_K) \\
V' &= X \setminus \overline{U'_K} = (X \setminus U'_K) \setminus C = \operatorname{Int}(X'_1),
\qquad X'_1 = X \setminus U'_K
\end{align*}
enfin $W = V \cap V' = \complement K$.

\begin{tikzcd}
U' \arrow[rr, no head, "C"] \arrow[dr, no head, "B'"'] & & U \arrow[dl, no head, "B"] \\
& W &
\end{tikzcd}

\note{à droite, une figure : un anneau autour d'un disque, le bord intérieur
coupé en $B'$, $C$, $B$, les régions marquées $U'$, $U$, $W$, $V$, $V'$ ;
une accolade désigne « $\operatorname{Int} K$ » ; non redessinée}

Alors on a : $U, U', W$ ouverts disjoints, $B, B', C$ fermés disjoints
— $U, U', W, B, B', C$ pseudo-partition de $X$
\[
\begin{cases}
\overline{U} \cap \overline{U'} = C, \quad \overline{U} \cap \overline{W} = B, \quad \overline{W} \cap \overline{U'} = B' \\
\dot{U}' = \overline{U'} \setminus U' = B' \cup C, \quad \dot{U} = B \cup C, \quad \dot{W} = B \cup B' \\
K = \overline{U \cup U'} = \complement W = (U \cup U') \cup B \cup B' \cup C, \quad
\operatorname{Int} K = U \cup U' \cup C
\end{cases}
\]
\note{la première égalité de la deuxième ligne commence par
$\overline{U^{*}} \cap \overline{W}$, lapsus probable}

4) \struck{\ill{}} Soient $U, U', W$ trois ouverts d'un espace $X$.
Conditions équivalentes :
\begin{enumerate}
\item[a)] $U, U', W$ mutuellement disjoints, réunion dense,
\struck{\ill{}} $\dot{U} \subset \overline{U' \cup W}$,
$\dot{U}' \subset \overline{U \cup W}$,
$\dot{W} \subset \overline{U \cup U'}$,
$\overline{U} \cap \overline{U'} \cap \overline{W} = \emptyset$,
\item[b)] Posant $C = \overline{U} \cap \overline{U'}$,
$B = \overline{U} \cap \overline{W}$, $B' = \overline{W} \cap \overline{U'}$,
les $U, U', W, B, B', C$ forment une pseudo-partition de $X$.
\end{enumerate}

On a alors les relations ci-dessus —

Je reviens \ill{} \ill{} données (\ill{} \ill{}) : $U, U', W$ satisfaisant
les conditions équ.\ a) b) — $K$, et $U_K, U'_K$ (fermé de $X$ ; ouverts de
$K$) satisfaisant
\[
\begin{cases}
\overline{U_K} = K \setminus U'_K, \quad \overline{\uncertain{U'_K}} = K \setminus U_K \\
C \cap \dot{K} = \emptyset \quad \text{car} \quad C = \overline{U_K} \cap \overline{U'_K} = K \setminus (U_K \cup U'_K)
\end{cases}
\]

\page{12}
\note{page traversée d'un long trait oblique ; le calcul du bas est en
outre barré de plusieurs traits}

\begin{tikzcd}
D(d) = \pi_0 \Sigma_d(\mathcal{X}) \arrow[r, "\alpha_d \sim"] \arrow[d, "\mathrm{P}_{d',d}"'] & E(d) \arrow[d, "r_{d',d}"] \\
D(d') \arrow[r, "\alpha_{d'} \sim"] & E(d')
\end{tikzcd}

\note{à gauche du diagramme, « $d \supset d'$ » écrit en colonne ;
l'étiquette de la flèche verticale de gauche est surchargée}

$E(d) \overset{\mathrm{df}}{=} \{e \in E \mid \varphi(M(e)) = d\}$

\[
D(\{i\}) = \pi_0 \Sigma(X_i, \partial X_i) \xrightarrow[\sim]{\ \alpha_i\ } E(\{i\}) = \varphi^{-1}(\{i\})
\]
avec, sous $\Sigma(X_i, \partial X_i)$ :
$\bigcup_{\mu\varphi(M(e)) \leq i} |E_e|$,
$\bigcup_{\mu\varphi(M(e)) < i} |E_e|$.

\[
\partial X_i = \bigcup_{\mu(\varphi(M(e))) < i} |E_e|
\]
\[
\partial X_i \overset{\mathrm{df}}{=} \bigcup_{j < i} X_j
= \bigcup_{j < i}\ \bigcup_{\substack{e \in E \\ \varphi M(e) \leq j}} |E_e|
= \bigcup_{\substack{e \in E \\ \exists j \text{ tel que } \varphi(M(e)) \leq j < i}} |E_e|
\]

\[
|K| = (K_0, K) \supset (L_0, L) = |L|, \qquad \Sigma(|K|, |L|) = |\Delta|
\]

\[
\begin{cases}
\tilde{K} \overset{\mathrm{df}}{=} \Sigma(K, L) = (\tilde{K}_0, \tilde{K}) \\
\tilde{K}_0 = \varinjlim_{\sigma \in K \setminus L} \sigma \ (\subset K_0) \xrightarrow{\ \pi\ } K_0 \\
\tilde{\Delta} = \operatorname{Im}\bigl((K \setminus L) \to \mathcal{P}(\Delta_0)\bigr)
\end{cases}
\]

\struck{$\ill{}\,|\tilde{K}| \xrightarrow{\ ?\ } \Sigma(|K|, |L|)$,
$= \varinjlim_{\sigma' \in \Delta} |\sigma'|$ ;
$\sigma' = \operatorname{Im}(\sigma)$, $\sigma \in K \setminus L$ ;
$|\sigma| \to |\sigma'| \to \Sigma$}
\note{la flèche porte un « $\sim$ » surmonté d'un « ? »}

\page{13}
\[
\Gamma_1 = \Gamma'_2 \cup \Gamma'_3, \qquad
\Gamma_2 = \Gamma'_3 \cup \Gamma'_1, \qquad
\Gamma_3 = \Gamma'_1 \cup \Gamma'_2
\]

a) $v_1, v_2, v_3$ \struck{appartiennent} contenus dans des régions
distinctes de $E - \Gamma$

$\Updownarrow$ ; a') $E - \Gamma$ a trois régions distinctes

b) $\Gamma_1, \Gamma_2, \Gamma_3$ sont des courbes de Jordan

$\Downarrow$

c) $\Gamma_1, \Gamma_2$ sont des courbes de Jordan

\note{la figure : deux grands arcs $\Gamma'_1$, $\Gamma'_2$ et l'arc
médian $\Gamma'_3$ se rejoignant en $x$ (en haut) et en bas ; un petit
cercle autour de $x$, découpé en trois secteurs $v_1, v_2, v_3$ ; non
redessinée}

a) $\Longleftrightarrow$ a')

b) $\Rightarrow$ c) \struck{$\Longleftrightarrow$} trivial

c) $\Rightarrow$ a)

\uncertain{Or} 2 pts de part et d'autre de $\Gamma'_3$ dans
\struck{\ill{} \ill{} $\Gamma'_3$ \ill{} dit} \ill{} \ill{} a fortiori
des régions distinctes de $E - \Gamma_1$ ($E - \Gamma_2$), \uncertain{donc}
a fortiori des régions distinctes de \struck{\ill{}}
$E - \Gamma_{\uncertain{2}}$ $= E - \Gamma$, donc $v_1$ et $v_2$ sont
des régions distinctes. De même, deux pts de part et d'autre de $\Gamma'_1$
($\Gamma'_2$) sont des régions distinctes de $E - \Gamma_2$ ($E - \Gamma_1$)
donc sont des régions distinctes de $E - \Gamma$, OK.

a) $\Rightarrow$ b) ; Prouvons que $\Gamma_3$ (p.\ ex.) est courbe de
Jordan.

\page{16}
\note{en haut à gauche, un croquis : deux points reliés par trois arcs
$\Gamma'_1$, $\Gamma'_3$, $\Gamma'_2$, les régions $V_2$, $V_1$ entre eux et
$V_3$ à l'extérieur, un petit cercle orienté autour de chaque point}

$\Gamma'_\alpha$, $\alpha \in \mathbb{I}$ \uncertain{ensemble} \ill{}
\ill{} : chemin simple joignant les deux points de $\partial = \{x, y\}$

$\Gamma_\alpha = \bigcup_{\beta \in A - \{\alpha\}} \Gamma'_\beta$ ;
\uncertain{déterminent} \ill{} \ill{} $B_\alpha$ de card.\ 2 (l'une des 2
origines, des pts $x$)

$B = \coprod_{i \in \mathbb{I}} B_i$ ; \struck{\ill{} \ill{} de 3
\ill{}}

Si $i \in I$ est donné, on lui associe $U_i$ la région de $\Gamma_i$ qui
contient $\Gamma'_i - \partial = \mathring{\Gamma}_i$, \uncertain{d'où}
\ill{} \ill{} de la \ill{}

\[
B \simeq I \times \{0, 1\}, \quad V_i :
\qquad
U_i = \mathring{\Gamma}_i \cup \bigcup_{j \in I \setminus \{i\}} U_j,
\qquad
V_i = \bigcap_{j \in I \setminus \{i\}} V_j
\]
\note{ces deux formules, telles qu'écrites, portent $U_j$ et $V_j$ sous les
signes de réunion et d'intersection}

$U_i, U_j, U_k$
\[
\left\{
\begin{array}{l}
V_1, V_2, V_3 \\
\Gamma^{\circ}_1, \Gamma^{\circ}_2, \Gamma^{\circ}_3 \\
\partial
\end{array}
\right.
\quad
\left\{
\begin{array}{l}
U_1 = V_2 \cup V_3 \cup \Gamma^{\circ}_1 \\
U_2 = V_3 \cup V_1 \cup \Gamma^{\circ}_2 \\
U_3 = V_1 \cup V_2 \cup \Gamma^{\circ}_3
\end{array}
\right.
\qquad
\begin{array}{l}
V_3 = U_1 \cap U_2 \\
V_2 = U_3 \cap U_1 \\
V_1 = U_1 \cap U_2
\end{array}
\quad
\left\{
\begin{array}{l}
\Gamma_1 = \Gamma^{\circ}_2 \cup \Gamma^{\circ}_3 \cup \partial \\
\Gamma_2 = \Gamma^{\circ}_3 \cup \Gamma^{\circ}_1 \cup \partial \\
\Gamma_3 = \Gamma^{\circ}_1 \cup \Gamma^{\circ}_2 \cup \partial
\end{array}
\right.
\]
\note{ainsi sur la page : $V_1 = U_1 \cap U_2$, comme $V_3$ ; on attend
$U_2 \cap U_3$}

$C$ courbe simple fermée dans surface connexe $E$. Conditions équ. :
\begin{enumerate}
\item[a)] $\operatorname{card} \pi_0(E - C) = 2$ ; (NB on a tjs
$\operatorname{card} \pi_0(E - C) \leq 2$)
\item[b)] $\exists$ ouvert $U \subset C$ tel que
$\dot{U}\ (\overset{\mathrm{df}}{=} \overline{U} - U) = C$ et que
\struck{\ill{}} $\overline{U} \neq E$ i.e.\
$V = U - \overline{U} \neq \emptyset$ i.e.\ $\overline{U} \neq E \setminus C$.
\item[c)] (étant donné un disque $D_x$ assez petit centré en $x \in C$), et
un point $x'$ et un point $x''$ $\in D_x - D_x \cap C$ de part et d'autre
de $D_x \cap C$…) $x'$ et $y$ n'appartiennent pas à la même composante
connexe de $X - C$, i.e.\ $\nexists$ arc dans $X - C$ joignant $x', y$.
(i.e.\ $\{x, y\} \to \pi_0(X - C)$ est surjectif…)
\end{enumerate}
\marginal{($U$ est alors une composante connexe de $X - C$), en regard de b)
et c)}
\note{« $U \subset C$ » et « $V = U - \overline{U}$ » sont ainsi sur la
page ; on attend $U \subset E$ et $V = E - \overline{U}$}

\page{18}
\emph{Proposition} Soient $\Gamma, \Gamma'$ deux J.-polygônes \uncertain{du}
plan $E$, tels que $\Gamma \cap \Gamma' = I$ soit un arc polygonal non
réduit à un point. Alors on est dans l'un des deux cas suivants
(s'excluant mutuellement)
\begin{enumerate}
\item[a)] $\Gamma - I \subset \operatorname{Ext}(\Gamma')$,
$\Gamma' - I \subset \operatorname{Ext}(\Gamma)$
 ($\Leftrightarrow$ $\operatorname{Int}\Gamma' \subset \operatorname{Ext}\Gamma$,
$\operatorname{Int}\Gamma \subset \operatorname{Ext}\Gamma'$)
\item[b)] $\Gamma - I \subset \operatorname{Ext}\Gamma'$,
$\Gamma' - I \subset \operatorname{Int}\Gamma$ (ou inversement)
($\Leftrightarrow$ $\operatorname{Int}\Gamma' \subset \operatorname{Int}\Gamma$,
$\operatorname{Ext}(\Gamma) \subset \operatorname{Ext}\Gamma'$ (ou
inv.))
\end{enumerate}

\struck{Dans les deux cas,} Soit
$\Gamma'' = (\Gamma \cup \Gamma') - (I - \partial I) = \overline{\Gamma - I} \cup \overline{\Gamma' - I}$,
qui est un polygône. C'est un J-polygône, et \ill{} :

\emph{Cas a)}
\begin{align*}
\operatorname{Int}\Gamma'' &= \operatorname{Int}\Gamma \cup \operatorname{Int}\Gamma' \cup (I - \partial I) \\
\operatorname{Ext}\Gamma'' &= \operatorname{Ext}\Gamma \cap \operatorname{Ext}\Gamma'
\end{align*}
[$E$ admet une partition en $\operatorname{Int}\Gamma$,
$\operatorname{Int}\Gamma'$, $\Gamma - I$, $\Gamma' - I$,
$I \setminus \partial I$, $\partial I$,
$(\operatorname{Ext}\Gamma \cap \operatorname{Ext}\Gamma')$]

\emph{Cas b)}
\begin{align*}
\operatorname{Int}\Gamma'' &= \operatorname{Int}\Gamma \setminus \operatorname{Int}\Gamma' \setminus (\Gamma' - I)
\quad \bigl[= \operatorname{Int}\Gamma \cap \operatorname{Ext}\Gamma'\bigr] \\
\operatorname{Ext}\Gamma'' &= \operatorname{Ext}\Gamma \cup \operatorname{Int}\Gamma' \cup (I \setminus \partial I)
\end{align*}
[$E$ admet une partition en
$\operatorname{Ext}\Gamma$, $(\operatorname{Int}\Gamma \cap \operatorname{Ext}\Gamma')$,
$\operatorname{Int}\Gamma'$,
\struck{\ill{}} $\Gamma - I$, $\Gamma' - I$, $I - \partial I$, $\partial I$]
\note{« $\operatorname{Ext}\Gamma$ » est écrit au-dessus de la parenthèse,
en ajout}

\page{20}
\note{Page « I » d'un mémoire d'une autre main, à l'encre bleue, signé et
daté du 15 juillet 1976, dont les pages XIV à XIX sont les pages 41 à 47 du
dossier ; il n'est pas recomposé ici. Titre : « Le théorème de Jordan et
l'existence de trois bons sommets pour un polygone de $\mathbb{R}^2$ ».
Théorème 1 (Jordan) : tout polygone est un J-polygone (déf.\ 1).
Théorème 2 (attribué avec un « ? » à Dehn, « non-publié ») : tout polygone
a au moins trois bons sommets, dont deux au moins à triangle intérieur.
§ 1, Préliminaires : un polygone est la réalisation géométrique dans
$\mathbb{R}^2$ d'un polygone abstrait connexe et non vide, de sommets
$S(\Gamma)$ et d'arêtes $A(\Gamma)$ ; un arc de polygone $\Gamma_{xy}$ est
la réalisation d'un 1-complexe connexe dont chaque sommet est incident à une
ou deux arêtes, exactement deux à une seule, ses extrémités $x, y$.
Remarque 1 : deux points distincts $x, y$ d'un polygone $\Gamma$ le coupent
en un seul couple d'arcs $\Gamma^1_{xy}, \Gamma^2_{xy}$ ; l'énoncé se
poursuit à la page suivante.}

\marginal{au crayon, d'une main non identifiée : « de $\mathbb{R}^2$ » du
titre est entouré, et « plan » écrit à la suite}

\marginal{à l'encre noire, de la main de ses feuillets de travail, en regard
de « d'un 1-complexe connexe » : « or.\ ? » ; un trait oblique suit
« connexe » dans le texte}

\marginal{à l'encre noire, de la même main, au-dessus de « connexe » dans la
définition du polygone : « et non vide », inséré ; au-dessus de « sous jacent
à $\Gamma$ » : « de $\mathbb{R}^2$ », inséré ; « aussi » et « couple »
soulignés}

\marginal{à l'encre noire, des points dans la marge gauche en regard de
plusieurs lignes du § 1 (« On entend par polygone », « Le terme arc de
polygone », « Par $\Gamma$ (resp.\ $\Gamma_{xy}$) », « Remarque 1 »)}

\marginal{au crayon, d'une main non identifiée : « Remarque 1 » est encadré
d'un trait tireté}

\marginal{à l'encre noire, au pied de la page, en regard de la fin de
« Remarque 1 » (« II) si… ») : « faux »}

\end{document}
